\documentclass[../../../include/open-logic-section]{subfiles}

\begin{document}

\olfileid{sth}{choice}{countablechoice}
\olsection{ગણનીય પસંદગી}

પસંદગી કે સુક્રમિતતાનો ઉપયોગ કર્યા વિના
આ સાબિત કરવું સરળ છે:

\begin{lem}[$\Zminus$માં]
દરેક સાન્ત ગણ માટે પસંદગી વિધેય છે.
\end{lem}

\begin{proof}
$a = \{b_1, \ldots, b_n\}$ લો. સરળતા માટે
દરેક $b_i \neq \emptyset$ ધારો. તેથી
$c_1 \in b_1, \ldots, c_n \in b_n$
એવી વસ્તુઓ $c_1, \ldots, c_n$ છે.
હવે \olref[z][pairs]{prop:pairsconsequences}
મુજબ $\{\langle b_1,
c_1\rangle , \ldots, \langle b_n, c_n\rangle\}$
ગણ છે; અને એ $a$ માટેનું પસંદગી વિધેય છે.
\end{proof}

પણ અનંત ગણોનો વિચાર કરતાં જ બાબત
વધુ અસ્પષ્ટ થાય છે. ઉદાહરણ તરીકે,
ઉપરના દાવાનું આ ``લઘુતમ'' વિસ્તરણ જુઓ:

\begin{defish}
\emph{ગણનીય પસંદગી.} દરેક \emph{ગણનીય}
ગણ માટે પસંદગી વિધેય છે.
\end{defish}

આ પસંદગીનો ખાસ કિસ્સો છે. જાણવા મળે છે
કે આ સિદ્ધાંતનો ઉપયોગ તેની સ્પષ્ટ જાણ
વિના ઘણી વાર થયો હતો. બે સરસ ઉદાહરણો
આ છે.\footnote{\citet[\S9.4]{Potter2004}
અને લૂકા ઇન્કુર્વાતી પાસેથી.}

\begin{ex}
એક સહજ વિચાર આ છે: કોઈ પણ ગણ $A$
માટે $\cardle{\omega}{A}$ હોય અથવા
કોઈ $n \in \omega$ માટે
$\cardeq{A}{n}$ હોય. દરેક ગણ સાન્ત કે
અનંત હોય છે એવા સહજ વિચારને કહેવાની
આ એક રીત છે. કૅન્ટર અને બીજા ઘણા
ગણિતજ્ઞોએ પુરાવા વિના આ દાવો કર્યો હતો.
સાવધાન રહીને આપણે
\olref[cardinals][classing]{generalinfinitycharacter}માં
તે સાબિત કર્યો. પરંતુ તે પુરાવામાં આપણે
$\ZFC$માં કામ કરતા હતા, કારણ કે કોઈ
પણ ગણ $A$ સુક્રમિત કરી શકાય છે અને
તેથી $\card{A}$ અસ્તિત્વમાં છે એમ માન્યું
હતું. એટલે આપણે પસંદગી સ્પષ્ટપણે માની હતી.

ખરેખર, \citet{Dedekind1888}એ આ દાવાનો
પોતાનો પુરાવો આ રીતે આપ્યો હતો:

\begin{thm}[$\Zminus + \text{ગણનીય પસંદગી}$માં]
કોઈ પણ $A$ માટે $\cardle{\omega}{A}$ હોય
અથવા કોઈ $n \in \omega$ માટે
$\cardeq{A}{n}$ હોય.
\end{thm}

\begin{proof}
દરેક $n \in \omega$ માટે
$\cardneq{A}{n}$ ધારો. તો ખાસ કરીને,
દરેક $n < \omega$ માટે ચોક્કસ
$\cardexpo{2}{n}$ ઘટકોવાળો ઉપગણ
$A_n \subseteq A$ છે. આ અનુક્રમ
$A_0, A_1, A_2, \ldots$ વાપરી દરેક
$n$ માટે આ વ્યાખ્યા આપીએ:
\[
	B_n = A_n \setminus \bigcup_{i < n} A_i.
\]
હવે આ નોંધો:
\begin{align*}
	\card{\bigcup_{i < n}A_i} 
	&\leq \card{A_0} + \card{A_1} + \ldots + \card{A_{n-1}}\\
	&=1 + 2 + \ldots + 2^{n-1}\\
	& = 2^n - 1\\
	& < 2^n = \card{A_n}
\end{align*}
તેથી દરેક $B_n$માં ઓછામાં ઓછો એક
ઘટક $c_n$ છે. વળી $B_n$ પરસ્પર
વિચ્છિન્ન છે; તેથી $c_n = c_m$ હોય
તો $n = m$. પણ દરેક $c_n \in A$.
આથી $f(n) = c_n$ એ
$\omega \to A$ અંતઃક્ષેપ છે.
\sourcecorrection{OLMD-149}{મૂળ અસમાનતાની
પહેલી પંક્તિમાં $\bigcup_{i<n}A_n$ છે;
યોગગણનો સૂચક $i$ હોવા છતાં તેનો
ઘટક $A_n$ લખાયો છે. આગળના
$A_0,\ldots,A_{n-1}$ સરવાળાને
સાચો આધાર આપવા $A_i$ કર્યું છે.}
\end{proof}
\noindent
ડેડેકિન્ડે ગણનીય પસંદગી વાપરી છે તે
જણાવ્યું નહોતું. પણ \emph{તમને} તેનો
ઉપયોગ દેખાયો? ફરી જુઓ. (ખરેખર:
\emph{ફરી જુઓ}.)

પુરાવાએ ગણનીય પસંદગી બે વાર વાપરી.
પ્રથમ, $A_0$, $A_1$, $A_2$, \dots\@ ગણોનો
અનુક્રમ મેળવવા; પછી દરેક~$B_n$માંથી
$c_n$ પસંદ કરવા. વળી પસંદગીનો આ
ઉપયોગ દૂર કરી શકાય તેમ નથી.
\citet[p.~138]{Cohen1966}એ બતાવ્યું કે
પસંદગીનું કોઈ પણ સ્વરૂપ ન હોય તો પરિણામ
નિષ્ફળ જાય છે. એટલે $\ZF$ સાથે એવા
ગણોનું અસ્તિત્વ સુસંગત છે જે
\emph{$\omega$ સાથે અતુલનીય} છે.
\end{ex}

\begin{ex}
\citeyear{Cantor1878}માં કૅન્ટરે કહ્યું
હતું કે ગણનીય ગણોનો ગણનીય યોગગણ
ગણનીય છે. તેમણે પુરાવો આપ્યો નહોતો;
કદાચ તેમને તે સ્પષ્ટ લાગતો હશે.
સાવધાન રહીને આપણે
\olref[card-arithmetic][simp]{kappaunionkappasize}માં
તેનું વધુ સામાન્ય સ્વરૂપ સાબિત કર્યું.
પરંતુ આપણા પુરાવામાં પસંદગી સ્પષ્ટપણે
માની હતી. અને ઓછા સામાન્ય દાવાના
પુરાવા માટે પણ ગણનીય પસંદગી જોઈએ.

\begin{thm}[$\Zminus + \text{ગણનીય પસંદગી}$માં]
દરેક $n \in \omega$ માટે $A_n$ ગણનીય
હોય, તો $\bigcup_{n <
\omega} A_n$ ગણનીય છે.
\end{thm}

\begin{proof}
સામાન્યતા ગુમાવ્યા વિના દરેક
$A_n \neq \emptyset$ ધારો. તેથી દરેક
$n \in \omega$ માટે !!a{surjection}
$f_n \colon \omega \to A_n$ છે.
$f(m, n) = f_n(m)$થી
$f \colon \omega \times \omega \to
\bigcup_{n < \omega} A_n$
વ્યાખ્યાયિત કરો. $\omega \times \omega$
ગણનીય છે
(\olref[sfr][siz][zigzag]{natsquaredenumerable})
અને $f$ !!a{surjection} છે,
તેથી પરિણામ મળે છે.
\end{proof}
\noindent
ગણનીય પસંદગીનો ઉપયોગ તમને દેખાયો?
$f_0$, $f_1$, $f_2$, \dots\ વિધેયોનો અનુક્રમ
પસંદ કરવા તે વપરાઈ.\footnote{પસંદગીનો
આવો જ ઉપયોગ
\olref[card-arithmetic][simp]{kappaunionkappasize}માં
થયો હતો, જ્યારે આપણે ``દરેક
$\beta \in \cardfont{a}$ માટે
!!a{injection} $f_\beta$ નિશ્ચિત લો''
એમ કહ્યું હતું.} અહીં પણ પસંદગીના
કોઈ પણ સિદ્ધાંત વિના પરિણામ નિષ્ફળ
જાય છે. ખાસ કરીને,
\citet{FefermanLevy1963}એ બતાવ્યું કે
ગણનીય ગણોનો ગણનીય યોગગણ
કદ~$\beth_1$ ધરાવે તે $\ZF$ સાથે
સુસંગત છે. પરંતુ રસેલનું આ વધુ રમૂજી
વર્ણન જુઓ:
\begin{quote}
  આ વાત એ કરોડપતિથી સ્પષ્ટ થાય છે, જે
  દરેક વાર બૂટની જોડી ખરીદે ત્યારે જ
  મોજાંની જોડી પણ ખરીદતો, અને બીજા
  કોઈ સમયે નહીં; બંને ખરીદવાનો તેને
  એટલો શોખ હતો કે અંતે તેની પાસે
  બૂટની $\aleph_0$ જોડીઓ અને મોજાંની
  $\aleph_0$ જોડીઓ થઈ ગઈ\dots\@
  બૂટમાં જમણો અને ડાબો જુદા ઓળખી શકીએ,
  તેથી દરેક જોડીમાંથી એક બૂટ પસંદ કરી
  શકીએ: બધા જમણા અથવા બધા ડાબા.
  પણ મોજાં માટે પસંદગીનો એવો કોઈ
  નિયમ સૂચાતો નથી. ગુણાકારની સ્વયંસિદ્ધિ
  [એટલે અસરકારક રીતે પસંદગી] ન માનીએ,
  તો દરેક જોડીમાંથી એક મોજું ધરાવતો
  કોઈ વર્ગ છે તેની ખાતરી નથી.
  \citep[p.~126]{Russell1919}
\end{quote}
ટૂંકમાં, આ સાબિત કરવા પસંદગીનું કોઈ
સ્વરૂપ જોઈએ: મોજાંની ગણનીય જોડીઓ હોય,
તો મોજાં પણ (માત્ર) ગણનીય હોય.
ખરેખર ગણનીય પસંદગી (અથવા સમકક્ષ
કંઈક) વિના ગણનીય ગણોનો ગણનીય
યોગગણ ગણનીય ન પણ હોઈ શકે.
\end{ex}

બોધ એ છે કે ગણનીય પસંદગી તેના
વપરાશકર્તાઓને બહુ જાણ થયા વિના
વારંવાર વપરાઈ. તત્ત્વચિંતનાત્મક પ્રશ્ન
આ છે: ગણનીય પસંદગીને \emph{વાજબી}
કેવી રીતે ઠરાવીએ?

સહજ સમર્થનનો એક પ્રયાસ અનંત પગલાંવાળા
કાર્યને અપીલ કરી શકે. પહેલી પસંદગી
$\nicefrac{1}{2}$ મિનિટમાં કરીએ, બીજી
$\nicefrac{1}{4}$ મિનિટમાં, \dots,
$n$મી પસંદગી $\nicefrac{1}{2^n}$
મિનિટમાં, \dots\@ તો $1$~મિનિટમાં
પસંદગીઓનો $\omega$-અનુક્રમ બનાવી
પસંદગી વિધેય વ્યાખ્યાયિત કરીશું.

પણ ખરેખર આવો વિચારપ્રયોગ આપણને
શું કહી શકે? પહેલાં તો તે ``પસંદ કરવું''
એ વિચારને બહુ શાબ્દિક રીતે લે છે.
વળી તે ગણિતને પરાભૌતિક શક્યતા
સાથે બાંધી દેતો લાગે છે.

વધુ મહત્ત્વનું એ છે કે તે માત્ર
ગણનીય પસંદગીને બદલે \emph{સમગ્ર}
પસંદગીને વાજબી ઠરાવશે નહીં. જો
\emph{દરેક} ગણ માટે પસંદગી વિધેય
જોઈએ, તો ``ગમે તેટલી ક્રમવાચક
લંબાઈનું અનંત પગલાંવાળું કાર્ય''
કરવું પડશે. સ્પષ્ટ કહીએ તો આ
વિચાર હાસ્યાસ્પદ છે.

\end{document}
